\documentclass{article}
\usepackage{iclr2026_conference,times}

\usepackage[utf8]{inputenc}
\usepackage[T1]{fontenc}
\usepackage{url}
\usepackage{booktabs}
\usepackage{amsfonts}
\usepackage{amsmath}
\usepackage{nicefrac}
\usepackage{microtype}
\usepackage{graphicx}
\usepackage{xcolor}
\usepackage{multirow}
\usepackage{subcaption}
\usepackage{wrapfig}
\usepackage{algorithm}
\usepackage{algpseudocode}
\usepackage{enumitem}
\usepackage{listings}
\usepackage[most]{tcolorbox}
\usepackage{hyperref}
\usepackage[noabbrev,capitalize]{cleveref}

% ---------------------------------------------------------------------------
% Colors
% ---------------------------------------------------------------------------
\definecolor{ebblue}{HTML}{1F5FA8}
\definecolor{ebgreen}{HTML}{2E7D32}
\definecolor{ebred}{HTML}{B5212F}
\definecolor{ebamber}{HTML}{B36B00}
\definecolor{ebgray}{HTML}{4A4A4A}
\definecolor{codebg}{HTML}{F5F5F5}
\definecolor{codeframe}{HTML}{D8D8D8}
\definecolor{keyword}{HTML}{0B5FB0}
\definecolor{string}{HTML}{2E7D32}
\definecolor{comment}{HTML}{8A8A8A}

\hypersetup{
  colorlinks=true,
  linkcolor=ebblue,
  citecolor=ebgreen,
  urlcolor=ebblue,
}

% ---------------------------------------------------------------------------
% Code listings: bash, yaml, python
% ---------------------------------------------------------------------------
\lstdefinelanguage{yaml}{
  keywords={true,false,null},
  keywordstyle=\color{keyword}\bfseries,
  basicstyle=\ttfamily\small,
  sensitive=false,
  comment=[l]{\#},
  morecomment=[s]{/*}{*/},
  commentstyle=\color{comment},
  stringstyle=\color{string},
  moredelim=[l][\color{ebblue}]{:},
}

\lstset{
  basicstyle=\ttfamily\footnotesize,
  backgroundcolor=\color{codebg},
  frame=single,
  rulecolor=\color{codeframe},
  framesep=5pt,
  xleftmargin=6pt,
  xrightmargin=2pt,
  breaklines=true,
  breakatwhitespace=true,
  showstringspaces=false,
  keywordstyle=\color{keyword}\bfseries,
  commentstyle=\color{comment}\itshape,
  stringstyle=\color{string},
  columns=fullflexible,
  keepspaces=true,
  postbreak=\mbox{\textcolor{codeframe}{$\hookrightarrow$}\space},
}

\lstdefinestyle{bash}{language=bash}
\lstdefinestyle{py}{language=Python}
\lstdefinestyle{yml}{language=yaml}

% Inline code shortcut
\newcommand{\code}[1]{\texttt{\small #1}}
\newcommand{\file}[1]{\texttt{\small #1}}

% ---------------------------------------------------------------------------
% Callout boxes
% ---------------------------------------------------------------------------
\newtcolorbox{tipbox}{
  enhanced, breakable, colback=ebgreen!5, colframe=ebgreen,
  boxrule=0.5pt, left=6pt, right=6pt, top=4pt, bottom=4pt,
  fonttitle=\bfseries, title={\faMagic\ Tip}, coltitle=ebgreen,
  attach title to upper, after title={\ \ }, titlerule=0pt,
}
\newtcolorbox{gotchabox}{
  enhanced, breakable, colback=ebred!5, colframe=ebred,
  boxrule=0.5pt, left=6pt, right=6pt, top=4pt, bottom=4pt,
  fonttitle=\bfseries, coltitle=ebred,
  attach title to upper, after title={\ \ }, titlerule=0pt,
  title={Gotcha.},
}
\newtcolorbox{notebox}{
  enhanced, breakable, colback=ebblue!4, colframe=ebblue,
  boxrule=0.5pt, left=6pt, right=6pt, top=4pt, bottom=4pt,
  fonttitle=\bfseries, coltitle=ebblue,
  attach title to upper, after title={\ \ }, titlerule=0pt,
  title={Note.},
}

% Track box: a project-track card with title argument
\newtcolorbox{trackbox}[1]{
  enhanced, breakable, colback=ebgray!3, colframe=ebgray,
  boxrule=0.7pt, left=7pt, right=7pt, top=5pt, bottom=5pt,
  fonttitle=\bfseries\color{white}, colbacktitle=ebgray,
  title={#1},
}

% Difficulty badges
\newcommand{\diffeasy}{\colorbox{ebgreen!18}{\textcolor{ebgreen}{\small\bfseries Warm-up}}}
\newcommand{\diffmed}{\colorbox{ebamber!18}{\textcolor{ebamber}{\small\bfseries Standard}}}
\newcommand{\diffhard}{\colorbox{ebred!15}{\textcolor{ebred}{\small\bfseries Ambitious}}}

% faMagic fallback (no fontawesome dependency)
\providecommand{\faMagic}{$\star$}

% ---------------------------------------------------------------------------
% Title
% ---------------------------------------------------------------------------
\title{EB-JEPA Hackathon Field Guide\\[0.3em]
\large A 24-Hour Sprint Through Joint-Embedding Predictive Architectures}

\author{Basile Terver and Théo Le Pendeven\\
EB-JEPA Hackathon\\
\texttt{github.com/facebookresearch/eb\_jepa}
}

\newcommand{\fix}{\marginpar{FIX}}
\newcommand{\new}{\marginpar{NEW}}

\iclrfinalcopy

\begin{document}

\maketitle
\fancyhead[L]{EB-JEPA Hackathon Field Guide}

\begin{abstract}
Welcome to the EB-JEPA Hackathon. Over the next 24 hours you will get hands-on
with Joint-Embedding Predictive Architectures (JEPAs): self-supervised models
that learn to predict in representation space rather than pixel space. This guide
is your field manual. It walks you through the three working examples shipped
with the \texttt{eb\_jepa} library (image representation learning, video
prediction, and action-conditioned world modeling with planning), shows you the
small set of components you will actually edit, and hands you a curated menu of
project tracks: most aimed at carrying the JEPA recipe to new continuous,
high-dimensional, noisy modalities, and some at pushing the vision and robotics
settings already in the codebase. It closes with the practical logistics:
compute policy on the shared B200 cluster, deliverables, and how projects are
judged. Read \cref{sec:gameplan} first, skim the rest, and start running code
within the hour.
\end{abstract}

% =============================================================================
\section{Welcome and the 24-Hour Game Plan}
\label{sec:gameplan}

This hackathon brings together roughly 100 participants in 25 teams of four,
sharing 72 NVIDIA B200 GPUs for 24 hours. The goal is not a leaderboard win; it
is to understand JEPAs by building with them. By the end you should be able to
explain, and have modified, the three core pieces of every JEPA: an
\emph{encoder} $f_\theta$ that maps observations to representations, a
\emph{predictor} $g_\phi$ that predicts future or alternate representations, and
a \emph{regularizer} $\mathcal{R}$ that stops those representations from
collapsing to a constant.

\paragraph{Why JEPAs, and why this codebase.}
Production JEPA stacks (I-JEPA, V-JEPA, V-JEPA~2)~\citep{IJEPA,VJEPA,VJEPA2} are
powerful but large and hard to navigate, and world-model implementations often
assume frozen pre-trained encoders and bespoke setups~\citep{DINO-WM,terver2026JEPAWMs}.
The \texttt{eb\_jepa} library~\citep{EBJEPA}
exists to close that gap: three self-contained examples, each trainable on a
single GPU in a few hours, that take you from static-image self-supervision to
temporal prediction to action-conditioned planning, all under one energy-based
objective (\cref{sec:recipe}). That makes it an ideal sandbox for a 24-hour
sprint: you can train a model end to end several times, ablate a component, and
actually see the effect before the clock runs out. By design \texttt{eb\_jepa} is
the prototyping rung of a ladder: once an idea works here, larger codebases pick
it up, \texttt{jepa-wms}~\citep{terver2026JEPAWMs} for planning with frozen
encoders on diverse benchmarks, and \texttt{stable-worldmodel}~\citep{stableworldmodel}
for a broad, controllable suite of world-model environments. Prototype small here;
scale and validate there.

\paragraph{What we want you to explore.}
The library ships with vision and robotics modalities. The headline challenge of
this hackathon is to take the \emph{JEPA recipe to data it was not built for}:
continuous, high-dimensional, noisy signals such as audio, physiological
recordings, climate and sensor fields, financial or scientific time series, and
point clouds. JEPAs are a natural fit here precisely because they predict in a
learned latent space and discard task-irrelevant noise, rather than reconstructing
every sample. Teams who prefer to stay in vision or robotics are very welcome to
push the existing examples with innovative ideas instead; \cref{sec:tracks}
offers a menu of both. Either way, the engineering surface is the same small set
of swappable components, and \cref{sec:extension} tells you exactly which files
to touch.

\paragraph{Scope realistically.}
Twenty-four hours is short. A successful project is one clear hypothesis,
tested with a clean before/after comparison, with a figure or number to show for
it. It is far better to get one new modality training without collapsing and to
\emph{understand why} than to half-wire three ambitious ideas. The single most
common failure mode in JEPA training is representation collapse; budget time for
it, watch the variance and prediction losses (\cref{sec:recipe}), and read the
gotchas in each example walkthrough.

\subsection{A suggested timeline}

\Cref{tab:timeline} is a sane default for a team of four, not a mandate. Adapt it
to your sleep schedule and your project's risk. The recurring theme: get real
code running in the first hour, lock your scope early, and stop \emph{building}
with several hours to spare so you can evaluate, visualize, and write up.

\begin{table}[h]
\centering
\caption{\textbf{A default 24-hour plan for a team of four.} Times are elapsed
hours from kickoff. ``Pods'' refers to splitting the team so that environment
setup, code reading, and a baseline run happen in parallel rather than in
series.}
\label{tab:timeline}
\small
\begin{tabular}{llp{8.4cm}}
\toprule
Phase & Hours & What to do \\
\midrule
Ignite      & 0--1   & Clone, install (\cref{sec:setup}), launch the \code{image\_jepa} smoke test on one GPU. Skim this guide. \\
Orient      & 1--3   & In parallel pods: (i) run all three examples; (ii) read the example you will build on; (iii) pick a track (\cref{sec:tracks}). \\
Lock scope  & 3--4   & Write a one-sentence hypothesis and a success metric. Sketch the data $\to$ encoder $\to$ predictor $\to$ loss path on paper. \\
Baseline    & 4--8   & Get \emph{something} training without collapsing: data loader yields the right tensor shape, loss decreases, variance stays alive. \\
Iterate     & 8--16  & One change at a time. Keep the previous run as the baseline; ablate; sweep two or three hyperparameters (\cref{sec:setup}). \\
Evaluate    & 16--21 & Lock the model. Run downstream eval / planning / probing. Make the figures. Re-run a seed for error bars if time allows. \\
Write up    & 21--24 & Push the repo (GitHub Classroom) and prepare the 10-minute presentation. State the hypothesis, the result, and one thing you learned about JEPAs. \\
\bottomrule
\end{tabular}
\end{table}

\begin{tipbox}
Treat the first training run as a plumbing test, not a science experiment. Use a
tiny model, a handful of epochs, and \code{logging.log\_wandb=false} to confirm the
pipeline runs end to end in minutes. Only then turn the knobs up. A pipeline that
runs at hour~5 beats a beautiful idea that first executes at hour~20.
\end{tipbox}

\subsection{How to read this guide}

\Cref{sec:recipe} is the two-page conceptual core: the unified JEPA objective and
the regularizers, enough theory to act. \Cref{sec:setup} gets you installed and
explains the launcher and the compute policy. \Cref{sec:examples} walks through
the three examples in the run-it / read-it / tweak-it pattern. \Cref{sec:extension}
is the reference you will return to most: the exact extension points for
encoders, predictors, regularizers, and datasets. \Cref{sec:tracks} is the
project menu. \Cref{sec:logistics} covers compute, deliverables, and judging.
For prior work, a curated per-track reading list lives in \code{references/README.md}
at the repo root (see \Cref{sec:tracks}). Throughout, \code{monospaced paths} point
at real files in the repository, and boxes flag \textbf{tips}, \textbf{gotchas},
and \textbf{notes} worth pausing on.
\section{The JEPA Recipe in Two Pages}
\label{sec:recipe}

This section is the whole conceptual toolkit you need to start. If you remember
one thing: a JEPA learns by \emph{predicting one representation from another},
and the only reason it does not cheat by mapping everything to a constant is the
regularizer. Everything else is architecture.

\begin{figure}[h]
\centering
\includegraphics[width=\linewidth]{figures/Archi_schema-eb-jepa.pdf}
\caption{\textbf{One architecture, three settings.} An encoder $f_\theta$ maps
inputs to representations; a predictor $g_\phi$ maps a context representation to a
predicted target $\hat{z}$; a cost $C$ measures prediction error in
representation space. \textbf{(a)} Image: two augmented views, $g_\phi$ is the
identity. \textbf{(b)} Video: $g_\phi$ predicts the next frame's representation.
\textbf{(c)} Action-conditioned video: the predictor is additionally conditioned
on an action $a_t$. \textbf{(d)} Planning: optimize an action sequence so the
imagined rollout reaches a goal representation $z_g$.}
\label{fig:archi}
\end{figure}

\subsection{The unified energy objective}

We view JEPAs through the lens of energy-based models~\citep{LeCunEBMTutorial2006}.
An energy function $E$ assigns a scalar to an input--output pair; low energy means
``compatible.'' For a JEPA, energy is simply prediction error in representation
space. Given an encoder $f_\theta$, a predictor $g_\phi$, and optional
conditioning $a$ with its own encoder $q_\omega$, the general energy is
\begin{equation}
E(x, x', a) = \mathcal{L}_{\text{pred}}\big(g_\phi(f_\theta(x), q_\omega(a)),\ f_\theta(x')\big),
\label{eq:energy}
\end{equation}
and training minimizes this energy plus a regularizer $\mathcal{R}$ that prevents
collapse:
\begin{equation}
\mathcal{L} = \mathcal{L}_{\text{pred}}\big(g_\phi(z, u), z'\big) + \lambda\,\mathcal{R}(z),
\qquad z = f_\theta(x),\ u = q_\omega(a).
\label{eq:general}
\end{equation}
The three examples are \cref{eq:general} with different choices of $g_\phi$ and
conditioning.

\paragraph{(a) Image-JEPA: invariance.} Two augmented views $x, x'$ of the same
image; the predictor is the identity, so the model just pulls the two
representations together,
\begin{equation}
\mathcal{L}_{\text{image}} = \|z - z'\|_2^2 + \lambda\,\mathcal{R}(z, z').
\label{eq:image}
\end{equation}
Low energy means the encoder has learned features invariant to the augmentation.

\paragraph{(b) Video-JEPA: temporal prediction.} The predictor sees a context of
$v{+}1$ frame representations and predicts the next one, summed over time,
\begin{equation}
\mathcal{L}_{\text{video}} = \sum_{t} \|g_\phi(z_{t-v:t}) - z_{t+1}\|_2^2 + \lambda\,\mathcal{R}(z_{1:T}).
\label{eq:video}
\end{equation}

\paragraph{(c) AC-video-JEPA: world modeling.} An action encoder produces
$u_t = q_\omega(a_{t-w:t})$ and the predictor is conditioned on it,
\begin{equation}
\mathcal{L}_{\text{world}} = \sum_{t} \|g_\phi(z_{t-v:t}, u_{t-v:t}) - z_{t+1}\|_2^2 + \lambda\,\mathcal{R}(z_{1:T}, u_{1:T}).
\label{eq:world}
\end{equation}
This is a latent dynamics model: given a state and an action, predict the next
state. It is what makes planning (\cref{sec:planning}) possible.

\begin{notebox}
In code this is one method. \file{eb\_jepa/jepa.py} defines a single
\code{JEPA} module whose \code{unroll(...)} call handles all three settings,
plus planning. Image-JEPA uses an identity predictor; video drops the action
encoder; AC-video uses everything. When you build a new modality you are picking
an encoder, a predictor, and a regularizer, then handing them to this same class.
\end{notebox}

\subsection{Preventing collapse: the regularizers}

If you only minimized prediction error, the encoder would map everything to a
constant: zero error, zero information. JEPAs in this library prevent that with
explicit regularization (rather than stop-gradient or EMA
tricks~\citep{BYOL,SimSiam}). Two families are implemented.

\paragraph{VICReg}~\citep{VICReg} uses two terms on a batch of embeddings
$Z \in \mathbb{R}^{N \times D}$. A \emph{variance} term keeps every feature
dimension from shrinking,
\begin{equation}
\mathcal{L}_{\text{var}}(Z) = \frac{1}{D}\sum_{j=1}^{D} \max\!\Big(0,\ \gamma - \sqrt{\operatorname{Var}(Z_{:,j}) + \epsilon}\Big),
\label{eq:var}
\end{equation}
with target standard deviation $\gamma$ (typically $1$); and a \emph{covariance}
term decorrelates dimensions so the model uses all its capacity,
\begin{equation}
\mathcal{L}_{\text{cov}}(Z) = \frac{1}{D(D-1)}\sum_{i\neq j} [C(Z)]_{i,j}^2,
\qquad C(Z) = \frac{1}{N-1}(Z - \bar{Z})^\top (Z - \bar{Z}).
\label{eq:cov}
\end{equation}
The full regularizer is $\mathcal{R}_{\text{VICReg}} = \alpha \mathcal{L}_{\text{var}} + \beta \mathcal{L}_{\text{cov}}$,
computed on a learned projection $r = h_\psi(z)$ rather than directly on $z$.

\paragraph{SIGReg}~\citep{balestriero2025lejepaprovablescalableselfsupervised}
takes a different route. It identifies the isotropic Gaussian $\mathcal{N}(0, I)$
as the optimal embedding distribution and enforces it by testing Gaussianity
along random 1D projections $\xi_p$,
\begin{equation}
\mathcal{R}_{\text{SIGReg}}(Z) = \frac{1}{P}\sum_{p=1}^{P} G(Z\xi_p),
\label{eq:sigreg}
\end{equation}
where $G$ is the Epps--Pulley Gaussianity statistic. It has a single
hyperparameter $\lambda$ and linear time and memory, which tends to make it more
forgiving to tune than VICReg's two coefficients (\cref{fig:image_sweep}).

\begin{gotchabox}
Watch the variance loss. If $\mathcal{L}_{\text{var}}$ shoots up and stays high
while the prediction loss drops to near zero, your encoder is collapsing: the
predictor is winning by making everything constant. The fix is almost always more
regularization weight or a healthier projector, not less. The image example
prints these per-term losses every epoch; the world-model example logs them too.
\end{gotchabox}

\subsection{Two ingredients specific to world models}

\paragraph{Multistep rollout.} Training a predictor only one step ahead, then
unrolling it autoregressively at test time, creates \emph{exposure bias}: the
model never saw its own predictions as input. The fix is to roll out $K$ steps
during training and sum the per-order losses,
\begin{equation}
\mathcal{L}_{\text{pred}} = \sum_{k=1}^{K} \mathcal{L}_k,
\qquad z^{(k)}_{t+1} = g_\phi\big(z^{(k-1)}_{t-v:t},\ u_{t-v:t}\big),\quad z^{(0)}_t = f_\theta(x_{t-w:t}),
\label{eq:rollout}
\end{equation}
where order $k$ counts predictor calls from a ground-truth representation. On
Moving MNIST, downstream Average Precision improves steeply with $K$ and plateaus
around $K{=}4$ (\cref{fig:ap}). In code, $K$ is the \code{model.steps} (video) or
\code{model.nsteps} (AC-video) config knob, consumed inside \code{JEPA.unroll}.

\paragraph{Extra world-model regularizers.} Action-conditioned training in
\emph{randomized} environments needs two more terms~\citep{sobal2022jepaslowfeatures}.
A temporal-similarity loss encourages smooth trajectories, and an inverse-dynamics
(IDM) loss~\citep{PathAMI,curiositydrivenexplorationselfsupervisedprediction}
forces the representation to retain action-relevant information by predicting the
action from consecutive states,
\begin{equation}
\mathcal{L}_{\text{sim}} = \sum_t \|z_t - z_{t+1}\|_2^2,
\qquad
\mathcal{L}_{\text{IDM}} = \sum_t \|a_t - \mathrm{MLP}(z_t, z_{t+1})\|_2^2.
\label{eq:idm}
\end{equation}
The full AC objective combines everything,
\begin{equation}
\mathcal{L} = \mathcal{L}_{\text{pred}} + \alpha\mathcal{L}_{\text{var}} + \beta\mathcal{L}_{\text{cov}} + \delta\mathcal{L}_{\text{sim}} + \omega\mathcal{L}_{\text{IDM}}.
\label{eq:full}
\end{equation}
The ablation in the paper is blunt about how load-bearing each term is: removing
the IDM loss collapses Two Rooms planning from $97\%$ to $1\%$ success, because
the encoder latches onto spurious background correlations instead of the agent.

\subsection{Planning as energy minimization}
\label{sec:planning}

Once you have a latent dynamics model, planning is optimization. To reach a goal
observation $x_g$, search for an action sequence whose imagined rollout
accumulates low energy against the goal representation,
\begin{equation}
E_{\text{plan}}(a_{0:H}; x_0, x_g) = \sum_{t=1}^{H} \|f_\theta(x_g) - \hat{z}_t\|_2^2,
\qquad \hat{z}_{t+1} = g_\phi(\hat{z}_{t-v:t}, u_{t-v:t}),\ \hat{z}_0 = f_\theta(x_0).
\label{eq:plan}
\end{equation}
Summing over \emph{all} intermediate steps (rather than only the final state)
rewards efficient paths and is robust to compounding prediction error; it beats
final-state-only cost by $8$ points on Two Rooms. The library solves
\cref{eq:plan} with two population-based optimizers, MPPI~\citep{MPPI} (default,
soft exponential weighting over elite trajectories) and CEM (hard elite
refitting), both in \file{eb\_jepa/planning.py}. You rarely write planning code;
you choose a planner config and a cost, and tune the horizon $H$, the number of
samples $N$, and the number of elites $Q$.
\section{Environment and Compute Setup}
\label{sec:setup}

Get this working in the first hour. The detailed compute policy for the shared
B200 cluster is in \cref{sec:logistics}; here we cover installation, the launch
workflow, and a smoke test.

\subsection{Install}

The library uses \href{https://docs.astral.sh/uv/}{\code{uv}} for package
management. The fastest path:
\begin{lstlisting}[style=bash]
git clone https://github.com/facebookresearch/eb_jepa.git
cd eb_jepa
uv sync                       # create .venv and install deps
source .venv/bin/activate
\end{lstlisting}
If you prefer conda for system packages:
\begin{lstlisting}[style=bash]
conda create -n eb_jepa python=3.12 -y && conda activate eb_jepa
uv pip install -e . --group dev   # editable install + pytest/black/isort
\end{lstlisting}
Then point the library at your data and checkpoint directories (add these to
\code{\textasciitilde/.bashrc}):
\begin{lstlisting}[style=bash]
export EBJEPA_DSETS=/path/to/datasets     # where datasets live / download to
export EBJEPA_CKPTS=/path/to/checkpoints  # where runs are written
\end{lstlisting}

\begin{tipbox}
Confirm the install with the test suite before you trust anything:
\code{uv run pytest tests/}. It exercises the loss equivalences, the planning
loop, and the JEPA output formats; if it passes, your environment is sane.
\end{tipbox}

\subsection{On the hackathon cluster (DALIA, B200)}
\label{sec:cluster}

\textbf{Repo + venv.} Work inside the shared project space \code{/lustre/work/pdl17890/}
--- \emph{not} your \code{\$HOME}, whose inode quota blocks \code{git}. The \code{uv}
virtualenv is \emph{architecture-specific} (login nodes are \code{x86\_64}, compute
nodes are \code{aarch64}), so rather than activating a single \code{.venv}, source the
provided \file{env.sh}: it selects an arch-specific \code{UV\_PROJECT\_ENVIRONMENT},
cache, and \code{uv} binary, plus \code{EBJEPA\_CKPTS}.
\begin{lstlisting}[style=bash]
cd /lustre/work/pdl17890/<you>/eb_jepa
source env.sh    # -> venvs/eb_jepa_{x86_64|aarch64}; then `uv pip install ...` as needed
\end{lstlisting}

\textbf{SLURM.} Submit with \code{--account=pdl17890 --partition=defq --gres=gpu:1};
exclude the frequently-OOM node with \code{--exclude=dalianvl03}. The cluster is
offline: set \code{WANDB\_DISABLED=true} \emph{and} put \code{log\_wandb: false} in
the config --- a CLI \code{--logging.log\_wandb=false} is parsed as the string
\code{"false"} (truthy) and still triggers a \code{wandb} error.

\textbf{Pre-downloaded datasets} (under the shared project space; group-readable):
\begin{lstlisting}[style=bash]
DSETS=/lustre/work/pdl17890/udl806719/datasets
$DSETS/the_well/gray_scott_reaction_diffusion  # The Well / gray_scott (Track 5)
$DSETS/speech_commands_v2                       # Speech Commands v2  (Track 3)
$DSETS/modelnet40/modelnet40_ply_hdf5_2048      # ModelNet40          (Track 10)
$DSETS/Neuro                                    # EEG corpora        (Tracks 1, 2, 6)
\end{lstlisting}
Point \code{EBJEPA\_DSETS} here (or read them directly); the maze/Two~Rooms data is
generated on the fly, so it has no path.

\subsection{Run something in two minutes}

Every example is a \code{python -m} module that reads a YAML config and accepts
dotted overrides on the command line. Start with the image example and a tiny,
no-logging configuration to prove the pipeline runs:
\begin{lstlisting}[style=bash]
python -m examples.image_jepa.main \
    --fname examples/image_jepa/cfgs/default.yaml \
    optim.epochs=2 logging.log_wandb=false
\end{lstlisting}
The three entry points are:
\begin{lstlisting}[style=bash]
python -m examples.image_jepa.main     --fname examples/image_jepa/cfgs/default.yaml
python -m examples.video_jepa.main     --fname examples/video_jepa/cfgs/default.yaml
python -m examples.ac_video_jepa.main  --fname examples/ac_video_jepa/cfgs/train/two_rooms/train.yaml
\end{lstlisting}

\begin{gotchabox}
All configs default to \code{logging.log\_wandb: true}. For local smoke tests set
\code{logging.log\_wandb=false} or you will hit a \code{wandb} login prompt. Also
note datasets download to \code{EBJEPA\_DSETS}; CIFAR-10 and Moving MNIST fetch
automatically on first run (Moving MNIST is $\sim$800\,MB and needs internet),
while Two Rooms is generated on the fly.
\end{gotchabox}

\subsection{The SLURM launcher}

For anything beyond a smoke test, submit to the cluster with the shared launcher
\file{examples/launch\_sbatch.py}. It wraps single jobs, three-seed sweeps, and
full hyperparameter grids.
\begin{lstlisting}[style=bash]
# One job (dev mode)
python -m examples.launch_sbatch --example image_jepa \
    --fname examples/image_jepa/cfgs/default.yaml --single

# Three seeds {1, 1000, 10000}, wandb-averaged (recommended default)
python -m examples.launch_sbatch --example image_jepa \
    --fname examples/image_jepa/cfgs/default.yaml --sweep my_experiment

# Full grid from the YAML's sweep.param_grid section
python -m examples.launch_sbatch --example image_jepa \
    --fname examples/image_jepa/cfgs/default.yaml --sweep my_grid \
    --full-sweep --use-wandb-sweep --array-parallelism 8
\end{lstlisting}
Swap \code{image\_jepa} for \code{video\_jepa} or \code{ac\_video\_jepa}. Key
flags: \code{--single} (one dev job), \code{--sweep NAME} (name your run group),
\code{--full-sweep} (read \code{sweep.param\_grid} from the config),
\code{--use-wandb-sweep} (the wandb sweep UI), and \code{--array-parallelism N}
(cap concurrent jobs, which matters a lot on a shared cluster, see
\cref{sec:logistics}).

\begin{gotchabox}
The launcher's default SLURM settings (partition, account, memory) at the top of
\file{examples/launch\_sbatch.py} point at the original authors' cluster. On the
hackathon cluster use \code{--account=pdl17890 --partition=defq} (and
\code{--exclude=dalianvl03}), as in \cref{sec:cluster}. Do not burn an hour
debugging a queue rejection; ask. See \cref{sec:logistics}.
\end{gotchabox}

\subsection{Where runs land, and how to compare them}

Runs are written under \code{\$EBJEPA\_CKPTS/\{example\}/} with an auto-named
experiment folder encoding the key hyperparameters, e.g.\
\code{resnet\_vicreg\_proj\_bs256\_ep300\_std1.0\_cov80.0\_seed1}. Three-seed
sweeps share a wandb run \emph{name} so you can group by name in the UI to get
mean and standard error automatically. To compare a before/after change, change
\emph{one} thing, keep the same seed set, and read the matched-epoch metric.

\begin{tipbox}
The recommended scientific unit is the three-seed sweep, not a single run. JEPA
training has real seed variance, and a one-point ``improvement'' from a single
run is usually noise. With B200s you can afford three seeds; budget for it.
\end{tipbox}
\section{The Three Examples: Run It, Read It, Tweak It}
\label{sec:examples}

The library is built around one idea: the training loop, the rollout, and the
losses are modality-agnostic and live in the shared \file{eb\_jepa/} package; each
example in \file{examples/} is a thin \file{main.py} that assembles an encoder, a
predictor, and a regularizer, then calls \code{JEPA.unroll}. Once you have read
one example, you have read them all. We go through them in order of increasing
complexity. For each: how to run it, the handful of files worth reading, the
config knobs that matter, and what to tweak first.

\begin{notebox}
There is no central ``registry'' of encoders or losses. Components are plain
\code{nn.Module}s wired together by hand inside each example's \file{main.py},
sometimes behind a small \code{if cfg.model.type == "..."} branch. To add your
own, you edit the builder in \file{main.py} (and occasionally add a class to
\file{eb\_jepa/architectures.py} or \file{eb\_jepa/losses.py}). \Cref{sec:extension}
is the map.
\end{notebox}

% ===========================================================================
\subsection{Image-JEPA: self-supervised features on CIFAR-10}
\label{sec:image}

\paragraph{What it teaches.} The purest form of the recipe: no predictor, no
time. Two augmented views of an image, pull their representations together, and
keep them from collapsing with VICReg or SIGReg. A linear probe trained online
reports classification accuracy each epoch ($\sim$91\% with SIGReg, $\sim$90\%
with VICReg at 300 epochs).

\paragraph{Run it.}
\begin{lstlisting}[style=bash]
# ResNet-18 + VICReg (default)
python -m examples.image_jepa.main --fname examples/image_jepa/cfgs/default.yaml
# ResNet-18 + SIGReg
python -m examples.image_jepa.main --fname examples/image_jepa/cfgs/sigreg.yaml
# ViT-S + VICReg
python -m examples.image_jepa.main --fname examples/image_jepa/cfgs/transformers.yaml
\end{lstlisting}
There is no separate eval step: the probe accuracy \code{val\_acc} is logged every
epoch during pretraining.

\paragraph{Read it.}
\begin{itemize}[leftmargin=1.4em,itemsep=1pt,topsep=2pt]
  \item \file{examples/image\_jepa/dataset.py} returns two augmented views per
    image (crop, color jitter, grayscale, solarize, flip).
  \item \file{examples/image\_jepa/main.py}: encoder built at lines
    \code{435--468}; loss chosen at \code{518--522}; the batch loop with the
    online probe at \code{275--321}; the epoch loop at \code{539--605}.
  \item \file{eb\_jepa/losses.py}: \code{VICRegLoss} and \code{BCS} (the SIGReg
    implementation) take two views and return a loss dict.
\end{itemize}

\paragraph{Key knobs (\code{cfgs/default.yaml}).}
\begin{lstlisting}[style=yml]
model:  {type: resnet, use_projector: true, proj_hidden_dim: 2048, proj_output_dim: 2048}
loss:   {type: vicreg, std_coeff: 1.0, cov_coeff: 80.0, lmbd: 10.0}  # lmbd used iff type=bcs
optim:  {epochs: 300, lr: 0.3, weight_decay: 1.0e-4}                  # LARS + cosine warmup
data:   {batch_size: 256}
\end{lstlisting}

\paragraph{Tweak it first.}
\begin{itemize}[leftmargin=1.4em,itemsep=1pt,topsep=2pt]
  \item Swap the regularizer: \code{loss.type=bcs loss.lmbd=10}. SIGReg has one
    coefficient and is more forgiving (\cref{fig:image_sweep}).
  \item Swap the backbone: \code{model.type=vit\_s}. Compare convergence and final
    accuracy at equal epochs.
  \item Ablate the projector: \code{model.use\_projector=false}. Expect a $\sim$3
    point accuracy drop, a clean demonstration of why the regularizer is computed
    in a projected space.
\end{itemize}

\begin{gotchabox}
Bad regularization weights collapse this model fast. The paper's sweep shows
\code{std=100, cov=100} dropping VICReg to $10\%$ (chance) while \code{std=1, cov=100}
reaches $90\%$. If \code{val\_acc} is stuck near $10\%$, you are collapsed, not
under-trained. Also: the ViT \code{patch\_size} is currently hardcoded in
\file{main.py}, so the \code{model.patch\_size} YAML key is ignored for ViT.
\end{gotchabox}

\begin{figure}[h]
\centering
\includegraphics[width=0.86\linewidth]{figures/aggregated_perf_image.pdf}
\caption{\textbf{Hyperparameter sensitivity on CIFAR-10.} SIGReg (right) stays
stable across a naive sweep while VICReg (left) reaches a similar peak but needs
more careful tuning. Collapsed runs flatline near chance accuracy. Use this as
your mental model when a new run will not improve: suspect collapse before
suspecting capacity.}
\label{fig:image_sweep}
\end{figure}

% ===========================================================================
\subsection{Video-JEPA: temporal prediction on Moving MNIST}
\label{sec:video}

\paragraph{What it teaches.} The predictor wakes up. Frames are encoded, the past
representations are context, and the model predicts the next frame's
representation, unrolled $K$ steps to fight exposure bias. A detection head and a
pixel decoder (both trained on detached features, purely for evaluation) let you
score Average Precision and watch the rollout.

\paragraph{Run it.}
\begin{lstlisting}[style=bash]
python -m examples.video_jepa.main --fname examples/video_jepa/cfgs/default.yaml
# fewer rollout steps, larger batch:
python -m examples.video_jepa.main --fname examples/video_jepa/cfgs/default.yaml \
    model.steps=2 data.batch_size=128
\end{lstlisting}

\paragraph{Read it.}
\begin{itemize}[leftmargin=1.4em,itemsep=1pt,topsep=2pt]
  \item \file{eb\_jepa/datasets/moving\_mnist.py}: \code{MovingMNISTDet}, two
    bouncing digits; the dataset, not the example, owns the data.
  \item \file{examples/video\_jepa/main.py}: model assembled at \code{128--142}
    (\code{ResNet5} encoder, \code{ResUNet} predictor wrapped by
    \code{StateOnlyPredictor}, \code{VCLoss}, \code{SquareLossSeq}); training step
    at \code{191--204}.
  \item \file{eb\_jepa/jepa.py}: \code{unroll} in \code{parallel} mode, the
    $K$-step loop, at \code{142--157}. This is the heart of temporal training.
\end{itemize}

\paragraph{Key knobs (\code{cfgs/default.yaml}).}
\begin{lstlisting}[style=yml]
model: {dobs: 1, henc: 32, dstc: 16, hpre: 32, steps: 4}   # steps = K rollout depth
loss:  {std_coeff: 10.0, cov_coeff: 100.0}
optim: {epochs: 50, lr: 1.0e-3}
data:  {batch_size: 64}
\end{lstlisting}

\paragraph{Tweak it first.}
\begin{itemize}[leftmargin=1.4em,itemsep=1pt,topsep=2pt]
  \item Sweep the rollout depth \code{model.steps} $\in \{1,2,4,8\}$ and plot AP
    versus prediction horizon (\cref{fig:ap}). This reproduces the paper's central
    video result in an afternoon.
  \item Replace the spatial \code{ResUNet} predictor with the GRU
    \code{RNNPredictor} (already in \file{eb\_jepa/architectures.py}) and compare.
\end{itemize}

\begin{figure}[h]
\centering
\includegraphics[width=0.62\linewidth]{figures/AP_vs_t.pdf}
\caption{\textbf{Multistep rollout helps, and saturates.} Average Precision of
the autoregressive prediction versus horizon, for models trained with $K \in
\{1,2,4,8\}$ rollout steps. Training with longer rollouts aligns training with
autoregressive inference; the Pareto knee is around $K{=}4$. A clean, cheap result
to reproduce or extend to a new temporal modality.}
\label{fig:ap}
\end{figure}

% ===========================================================================
\subsection{AC-Video-JEPA: a world model you can plan with}
\label{sec:acvideo}

\paragraph{What it teaches.} The full stack: an action-conditioned latent dynamics
model, the two world-model regularizers ($\mathcal{L}_{\text{sim}}$ and the
critical $\mathcal{L}_{\text{IDM}}$), and goal-conditioned planning by energy
minimization. The agent navigates a ``Two Rooms'' environment whose wall and door
are randomized per trajectory, so the task is non-monotonic (sometimes you must
move \emph{away} from the goal to reach it). The best model plans at $97\%$ success
with MPPI.

\paragraph{Run it (train, then plan).}
\begin{lstlisting}[style=bash]
# Train the world model
python -m examples.ac_video_jepa.main --fname examples/ac_video_jepa/cfgs/train/two_rooms/train.yaml
# Evaluate planning on a trained checkpoint (runs unroll + planning eval)
python -m examples.ac_video_jepa.main \
    --meta.model_folder /path/to/run --meta.eval_only_mode True
\end{lstlisting}
Planning uses MPPI by default (\code{cfgs/planning/two\_rooms/planning\_mppi.yaml}); point
\code{eval.plan\_cfg\_path} at \code{cfgs/planning/two\_rooms/planning\_cem.yaml} for CEM.

\paragraph{Read it.}
\begin{itemize}[leftmargin=1.4em,itemsep=1pt,topsep=2pt]
  \item \file{eb\_jepa/datasets/two\_rooms/}: the environment (\code{DotWall}), the
    dataset, and the data-generation pipelines (online / streaming / offline).
  \item \file{examples/ac\_video\_jepa/main.py}: model built at \code{176--230}
    (\code{ImpalaEncoder}, \code{RNNPredictor}, action encoder
    \code{nn.Identity()}, \code{InverseDynamicsModel}, \code{VC\_IDM\_Sim\_Regularizer});
    training step calling \code{jepa.unroll(..., unroll\_mode="autoregressive")} at
    \code{334--343}.
  \item \file{eb\_jepa/planning.py}: \code{CEMPlanner} and \code{MPPIPlanner}, the
    goal-distance objective \code{ReprTargetDistMPCObjective}, and the eval entry
    points \code{main\_eval} / \code{main\_unroll\_eval}.
\end{itemize}

\paragraph{Key knobs.}
\begin{lstlisting}[style=yml]
# cfgs/train/two_rooms/train.yaml
model:  {encoder_architecture: impala, dobs: 2, dstc: 32, nsteps: 8}  # nsteps = K
  regularizer: {cov_coeff: 8, std_coeff: 16, sim_coeff_t: 12, idm_coeff: 1}
optim:  {epochs: 12, lr: 0.001}
data:   {batch_size: 384}
# cfgs/planning/two_rooms/planning_mppi.yaml
planner: {planner_name: mppi, plan_length: 90, n_iters: 20, num_samples: 200, num_elites: 20, temperature: 0.005}
\end{lstlisting}

\paragraph{Tweak it first.}
\begin{itemize}[leftmargin=1.4em,itemsep=1pt,topsep=2pt]
  \item Reproduce the headline ablation: set \code{model.regularizer.idm\_coeff=0}
    and watch planning success collapse toward $1\%$. Then zero
    \code{std\_coeff}, \code{cov\_coeff}, \code{sim\_coeff\_t} one at a time.
  \item Compare planners (MPPI vs CEM) and the cost design
    (\code{planning\_objective.sum\_all\_diffs} true vs false, i.e.\ cumulative vs
    final-state cost). Cumulative is worth $\sim$8 points.
  \item Trade planning compute for success: sweep \code{num\_samples} and
    \code{plan\_length} and report success-versus-wall-clock.
\end{itemize}

\begin{gotchabox}
The action dimension \code{action\_dim=2} is hardcoded in several places
(\file{main.py}, \code{InverseDynamicsModel}, the planners). If you retarget this
example to a new control problem, grep for \code{action\_dim} and change all of
them together. And remember: in randomized environments the IDM loss is what keeps
the encoder honest, do not drop it.
\end{gotchabox}
\section{Extension Points: The Files You Will Actually Edit}
\label{sec:extension}

This is the reference you will come back to. Extending the library is not
mysterious: there is no plugin system to learn. You add a plain \code{nn.Module}
and wire it into an example's \file{main.py} builder. \Cref{tab:extension} is the
``I want to\ldots{} $\to$ edit this'' map; the rest of the section explains the
two contracts that matter (the encoder/predictor I/O shape, and the regularizer
return signature).

\begin{table}[h]
\centering
\caption{\textbf{Where to make each kind of change.} ``Builder'' means the model
assembly block near the top of the relevant \file{examples/*/main.py}.}
\label{tab:extension}
\small
\begin{tabular}{p{4.6cm}p{9.2cm}}
\toprule
I want to\ldots & Edit \\
\midrule
Add an encoder for a new modality & New \code{nn.Module} in \file{eb\_jepa/architectures.py}; wire it into the builder (\file{image\_jepa/main.py}$\sim$435, or the hardcoded constructor in \file{video\_jepa/main.py}$\sim$130 / \file{ac\_video\_jepa/main.py}$\sim$186). \\
Add a predictor & New \code{nn.Module} in \file{architectures.py} exposing \code{is\_rnn} and \code{context\_length}; wire into the builder. \\
Add / change a regularizer & New class in \file{eb\_jepa/losses.py}; instantiate in the builder, or add an \code{elif cfg.loss.type} branch (image). \\
Add a dataset / modality & New \code{Dataset} (or \code{TrajDataset}) under \file{eb\_jepa/datasets/}; for the world-model path, add an \code{env\_name} branch in \file{eb\_jepa/datasets/utils.py:init\_data} (it currently rejects anything but \code{two\_rooms}); for the image path, edit the dataset block in \file{image\_jepa/main.py}$\sim$398. \\
Add a planner or planning cost & New class in \file{eb\_jepa/planning.py} plus an entry in \code{planner\_name\_map} / \code{objective\_name\_map}. \\
Change rollout depth, coefficients & Config only: \code{model.steps}/\code{model.nsteps}, \code{std\_coeff}, \code{cov\_coeff}, \code{sim\_coeff\_t}, \code{idm\_coeff}, \code{loss.lmbd}. \\
\bottomrule
\end{tabular}
\end{table}

\subsection{The central object: \texttt{JEPA.unroll}}

Everything funnels through one class, \code{JEPA} in \file{eb\_jepa/jepa.py},
constructed as \code{JEPA(encoder, action\_encoder, predictor, regularizer,
predcost)}. Its \code{unroll(observations, actions, nsteps, unroll\_mode, ...)}
method does, in order: encode the observations to \code{state}; if computing loss,
call \code{regularizer(state, actions)}; encode the actions; then roll the
predictor forward and accumulate the prediction cost. Two rollout modes:

\begin{itemize}[leftmargin=1.4em,itemsep=2pt,topsep=2pt]
  \item \code{unroll\_mode="parallel"} (used by \code{video\_jepa}): convolutional
    predictors predict all timesteps at once, refeeding a few ground-truth context
    frames; the $K$-step loop is at \file{jepa.py}\code{:142--157}.
  \item \code{unroll\_mode="autoregressive"} (used by \code{ac\_video\_jepa} and
    any RNN predictor): step-by-step with a sliding context window; loop at
    \file{jepa.py}\code{:161--190}.
\end{itemize}
\code{JEPA} chooses behavior by reading \code{getattr(predictor, "is\_rnn", False)}
and \code{predictor.context\_length}. So your predictor mostly needs to declare
those two attributes correctly and the rest follows.

\subsection{Contract 1: encoder and predictor tensor shapes}

The temporal examples speak in 5D tensors \code{[B, C, T, H, W]} (batch, channels,
time, height, width). The mixin \code{TemporalBatchMixin} in
\file{eb\_jepa/nn\_utils.py} folds time into the batch for you, so a 2D encoder can
be written as if it only ever sees \code{[B, C, H, W]}; subclass it and implement
\code{\_forward}. Encoders output a representation \code{[B, D, T, H', W']} (the
\code{ImpalaEncoder} pools to \code{H'{=}W'{=}1}, i.e.\ one vector per frame). A
predictor implements \code{forward(state, action) -> [B, D, T, ...]} and declares:

\begin{lstlisting}[style=py]
class MyPredictor(nn.Module):
    is_rnn = False          # True -> autoregressive single-step unroll
    context_length = 2      # how many context frames it consumes
    def forward(self, state, action):
        ...                 # return predicted next-step representation(s)
\end{lstlisting}

\begin{tipbox}
For a new modality with no spatial grid (a 1D time series, a feature vector per
step), the cleanest path is the GRU \code{RNNPredictor} (it sets
\code{is\_rnn=True}, \code{context\_length=0}) plus an encoder that outputs
\code{[B, D, T, 1, 1]}. You inherit the autoregressive rollout and planning for
free. Reshape your samples to that 5D convention and most of the spatial machinery
becomes a no-op.
\end{tipbox}

\subsection{Contract 2: the regularizer return signature}

There are two regularizer conventions, matching the two settings:

\begin{itemize}[leftmargin=1.4em,itemsep=2pt,topsep=2pt]
  \item \textbf{Sequence regularizers} (video / world model) accept
    \code{(state, actions)} and return a triple
    \code{(weighted\_loss, unweighted\_loss, loss\_dict)}. This is what
    \code{JEPA.unroll} expects. Models: \code{VCLoss} (variance+covariance) and
    \code{VC\_IDM\_Sim\_Regularizer} (the full world-model regularizer).
  \item \textbf{Two-view SSL losses} (image) accept \code{(z1, z2)} and return a
    dict with a \code{"loss"} key plus per-term entries. Models: \code{VICRegLoss}
    and \code{BCS} (SIGReg). The image training loop logs every extra key in the
    dict automatically, so new sub-losses show up in logging with no extra work.
\end{itemize}

Match whichever contract fits your setting and the rest of the loop is untouched.
Building blocks you can reuse to compose a new regularizer live in
\file{eb\_jepa/losses.py}: \code{HingeStdLoss} (variance), \code{CovarianceLoss},
\code{TemporalSimilarityLoss}, \code{InverseDynamicsLoss}, and the \code{BCS}
Gaussianity statistic.

\subsection{Adding a dataset: the one real dispatch}

For image-style two-view data, edit the dataset block in
\file{image\_jepa/main.py} (around line 398) to branch on \code{cfg.data.dataset}
and adjust \code{num\_classes} and the normalization constants. For the
world-model path, \file{eb\_jepa/datasets/utils.py:init\_data} is the single point
that maps an \code{env\_name} to data loaders, and it currently raises for any
\code{env\_name} other than \code{two\_rooms}. Adding a new environment or
trajectory dataset means implementing a \code{TrajDataset} (which yields
\code{(obs, action, state, ...)} slices via \code{TrajSlicerDataset}) and adding
your branch there. This is the most involved extension, so scope it early.

\begin{gotchabox}
Per-feature normalization is not optional for new modalities. The variance term in
VICReg measures per-dimension standard deviation; if one input channel dwarfs the
others (common in EEG, audio, sensor fusion), it will dominate both the encoder and
the regularizer. Normalize each channel/feature (z-score or robust scaling) before
it enters the encoder.
\end{gotchabox}
\section{Project Tracks}
\label{sec:tracks}

Below is a menu, not a syllabus. Pick one track, or use one as a springboard for
your own idea. Each track names a difficulty, the example it forks, a 24-hour
goal, the files you will touch, and a success metric. Difficulty badges:
\diffeasy{} is reachable by most teams and a safe first project; \diffmed{} is a
solid full-day project; \diffhard{} is high-risk, high-reward, attempt it only if
your baseline is already running.

\begin{notebox}
\textbf{The reading list.} A curated, per-track bibliography lives in
\code{references/README.md} at the repo root: one folder per paper
(\code{references/paper/<slug>/} with the PDF and a structured \code{SUMMARY.md}),
organized into a \textbf{Core} set (EB-JEPA and the JEPA world-model line it builds
on) plus three maturity tiers, every row tagged with the modality or track it
serves. Find the prior work for your track there, and read its \code{SUMMARY.md}
before the PDF.
\end{notebox}

\paragraph{What makes a good new-modality target.} The headline theme is
\emph{continuous, high-dimensional, noisy} data, exactly where predicting in pixel
or sample space is wasteful and predicting in a learned latent space shines. A
good target has: (i) abundant unlabeled data for pretraining; (ii) a small labeled
set for a downstream probe; (iii) temporal or multi-view structure the predictor
can exploit; and (iv) enough noise that reconstruction would be a bad idea. EEG
(Tracks~1--2) are the worked example of all four.

% ===========================================================================
\subsection{Group A: Carry the JEPA recipe to a new modality}

\begin{trackbox}{Track 1 (Flagship) --- Read the Brainwaves: abnormality \& epilepsy detection \hfill \diffeasy{}/\diffmed{}}
\textbf{Usefulness in neuroscience.} This track uses \textbf{electroencephalography
(EEG)}: non-invasive \emph{neurophysiological recordings} of the brain's electrical
activity, picked up by a few dozen electrodes on the scalp. Millions of clinical EEGs
are read by hand every year; a model that recovers the clinical label from the raw
signal with no hand-engineered features is a step toward automated triage.

\smallskip
\textbf{The idea.} Pretrain a JEPA on raw clinical EEG \emph{without labels}, freeze
it, and linear-probe whether the embeddings carry a clinical label that applies to the
\emph{whole recording}: normal vs abnormal (TUAB) and epilepsy vs not (TUEP). This mirrors the EEG
foundation model LaBraM~\citep{LaBraM} but swaps its reconstruction objective --- a
vector-quantized tokenizer that rebuilds the Fourier spectrum --- for a pure
joint-embedding one: \emph{predict patch representations in latent space}, never the raw
signal. That is the point of a JEPA: predicting in representation space lets the target
encoder discard the unpredictable parts of low signal-to-noise EEG, rather than forcing a
(spectral) reconstruction of it.

\smallskip
\textbf{Data (preprocessed, ready to probe).} Two labeled corpora, one label per
recording, already preprocessed on the cluster (each folder ships a
\code{<NAME>\_README.md} with the exact task, label layout, and preprocessing):
\begin{itemize}[leftmargin=1.4em,itemsep=1pt,topsep=2pt]
  \item \textbf{TUAB} (\code{TUAB\_PREPROCESSED}) --- binary normal/abnormal,
    2{,}717/276 train/eval recordings, \emph{official patient-disjoint split}.
  \item \textbf{TUEP} (\code{TUEP\_PREPROCESSED}) --- patient-level epilepsy vs
    no-epilepsy, 200 patients (100/100), a \emph{weak} label: the diagnosis is per
    patient, but epileptiform activity is intermittent, so most windows (often whole
    recordings) of an epileptic patient carry no visible sign of it. You build a
    subject-disjoint split and pool per-window scores into a patient decision.
\end{itemize}
Both: 19-channel 10--20 montage, 200\,Hz, $0.1$--$75$\,Hz bandpass, 60\,Hz notch
(TUEP average-referenced). Only \emph{windowing} and \emph{per-channel
normalization} are left to you. You pretrain the JEPA on this same EEG while
\emph{ignoring} the labels; the labels come back only for the linear probe.

\smallskip
\textbf{Two framings (start with the first).}
\begin{itemize}[leftmargin=1.4em,itemsep=1pt,topsep=2pt]
  \item \emph{Temporal} (video-JEPA style, \cref{sec:video}): cut each channel into
    1-second patches; encode the context window; predict the \emph{next} 1-second
    window's representation across all channels. Next-frame prediction with
    frames $=$ time windows.
  \item \emph{Augmented-view} (image-JEPA style, \cref{sec:image}): take an epoch
    (say 10\,s) as one ``image'' and form \emph{two augmented views} of it (two time
    crops, optionally $+$ masking, amplitude scaling, channel dropout, Gaussian noise);
    predict one view's representation from the other. The augmentation set is the main
    design knob.
\end{itemize}

\smallskip
\textbf{Files.} A new \code{Dataset} under \file{eb\_jepa/datasets/} that loads the
preprocessed EDF (\code{MNE-Python}) and yields one epoch per sample in eb\_jepa's
\code{[B, C, T, 1, 1]} layout --- \code{B} batch, \code{C} EEG channels, \code{T} time
samples, the trailing \code{1, 1} the singleton spatial dimensions of the
\code{[B, C, T, H, W]} convention (EEG has no spatial grid). Then a small conv-patch or
1D-conv encoder in \file{architectures.py}; the GRU \code{RNNPredictor} for the temporal
framing; a linear probe modeled on \file{examples/image\_jepa/eval.py}.

\smallskip
\textbf{Metric.} Balanced Accuracy / AUROC at the \textbf{window (epoch) level} --- each
window inherits its recording's label and is scored individually, as in recent EEG
foundation models such as LaBraM and BIOT (TUAB on its official patient-disjoint eval;
TUEP on your subject-disjoint split). Context: on other models, a supervised CNN-Transformer reaches $\sim$0.78 Balanced
Accuracy on TUAB, BIOT $\sim$0.80, and fine-tuned LaBraM-Base $\sim$0.81; a frozen-JEPA
\emph{linear probe} clearing $\sim$0.78 is a strong result.

\smallskip
\textbf{Benchmark (optional).} Put the frozen-JEPA probe next to a foundation model: run
the same linear probe on a pretrained \textbf{LaBraM}~\citep{LaBraM} (fine-tuned) over
your split and compare. A frozen JEPA matching a fine-tuned foundation model could be a
headline result.

\smallskip
\textbf{Gotchas.} \textbf{Evaluate only on held-out patients} (\emph{patient-disjoint}):
EEG carries a per-subject ``fingerprint'', so an overlapping split decodes \emph{who} not
\emph{what}; keep test patients out of both pretraining and the probe fit. The recordings
ship continuous --- \emph{epoch} them into fixed-length windows (e.g.\ 10\,s) before
pretraining or probing. Per-channel normalization is mandatory; use a class-aware metric
(eval is mildly imbalanced); score \emph{per window} --- each window carries its
recording's label, so report metrics over windows, not over collapsed recordings. For
TUEP the label is \emph{weak} --- most windows of an epileptic patient look normal, so
per-window scores are noisy and how you optionally pool them into a patient-level call
is the interesting part.

\smallskip
\textbf{poc:} TUAB frozen linear-probe \textbf{AUROC 0.877} (patient-disjoint split,
near supervised SOTA). TUEP (patient-level epilepsy) is preprocessed and ready as a
second, independent task --- its own encoder, pretrained and probed within TUEP
(\emph{IID}: pretraining and evaluation on the same corpus).
\emph{Data:} \code{\$DSETS/Neuro/\{TUAB,TUEP\}\_PREPROCESSED} (each with its
\code{*\_README.md}).

\smallskip
\textbf{Stretch (cross-dataset transfer).} Beyond the IID protocol above, the two
corpora enable a real generality test: pretrain the JEPA on one (say TUAB), freeze it,
then train a \emph{fresh} linear probe on the other's labels (the label spaces differ,
so only the encoder transfers). Does a representation learned on one cohort generalize
to another? That transfer is exactly what a foundation model claims.
\end{trackbox}

\begin{trackbox}{Track 2 --- When the Brain Misfires: artifact, epileptiform \& seizure detection \hfill \diffmed{}/\diffhard{}}
\textbf{Usefulness in neuroscience.} Like Track~1 this works on \textbf{EEG}
(electroencephalography --- scalp recordings of the brain's electrical activity), but the
clinically meaningful content of a recording is mostly \emph{events} --- spikes,
seizures, artifacts --- sparse in time and localized to a few channels; finding them is
the daily work of an epileptologist.

\smallskip
\textbf{The idea (analogous to Track~1 --- pickable on its own).} Same recipe: pretrain a
JEPA on raw EEG \emph{without labels}, freeze it, and linear-probe whether the embeddings
carry the label --- a pure joint-embedding objective (\emph{predict patch representations
in latent space}, never the raw signal), the same structural alternative to the
reconstruction tokenizer of LaBraM~\citep{LaBraM}. Pretrain with either framing from Track~1 (temporal next-window
or augmented-view). What changes from Track~1 is \emph{where} the label lives: not
one label for the whole recording, but \emph{event-level} labels sitting inside the
recording on \emph{specific channels} (see the annotation gotcha below). Three corpora of
rising richness:
\begin{itemize}[leftmargin=1.4em,itemsep=1pt,topsep=2pt]
  \item \textbf{TUAR} (\code{TUAR\_PREPROCESSED}) --- per-segment \emph{artifact}
    detection/typing (eye, muscle, chew, shiver, electrode), 290 recordings; gross
    and visible, a noise-robust warm-up into the event setting.
  \item \textbf{TUEV} (\code{TUEV\_PREPROCESSED}) --- 6-class event classification
    (spike/sharp-wave, two periodic-discharge types GPED/PLED, eye-movement, artifact,
    background), 359/159 train/eval, imbalanced.
  \item \textbf{TUSZ} (\code{TUSZ\_PREPROCESSED}) --- \emph{seizures}, the richest:
    each recording ships \emph{two} annotations, a binary whole-brain
    \code{seiz}/\code{bckg} timeline (detection) and a per-channel \emph{seizure
    type} (8 classes, typing); 7{,}515 recordings, official train/dev/eval,
    $\sim$13\% seizure-bearing.
\end{itemize}

\smallskip
\textbf{The annotation gotcha (read the README).} TUEV/TUAR/TUSZ events are
time-stamped on a \emph{bipolar} montage: each annotation gives a channel that names a
\emph{pair} of electrodes (not a row of the signal), a start and end time, and a label
--- so turning those annotations into training samples is itself part of the task. Each
folder's \code{<NAME>\_README.md} spells this out.

\smallskip
\textbf{Data (preprocessed).} All three at 200\,Hz, $0.1$--$75$\,Hz bandpass, 60\,Hz
notch (TUSZ average-referenced; 19--23 channels depending on the bipolar montage the
labels need). The recordings ship continuous --- you must first \emph{epoch} them into
fixed-length windows; windowing and per-channel normalization are left to you.

\smallskip
\textbf{Gotcha (patient-disjoint).} Evaluate only on held-out patients --- EEG carries a
per-subject ``fingerprint'', so an overlapping split decodes \emph{who} not \emph{what};
keep test patients out of pretraining and the probe fit. TUEV/TUSZ ship official splits;
build your own (patient-disjoint) for TUAR.

\smallskip
\textbf{Files.} A \code{Dataset} that loads the preprocessed EDF (\code{MNE-Python}) and
yields one epoch per sample in eb\_jepa's \code{[B, C, T, 1, 1]} layout (\code{B} batch,
\code{C} EEG channels, \code{T} time samples; the trailing \code{1, 1} are the singleton
spatial dims of the \code{[B, C, T, H, W]} convention). Then a conv-patch / 1D-conv
encoder in \file{architectures.py}; the GRU \code{RNNPredictor} for the temporal framing;
a linear probe modeled on \file{examples/image\_jepa/eval.py} --- \emph{plus} per-event
window extraction from the \code{.csv} / \code{.csv\_bi} annotations (the event-specific
part).

\smallskip
\textbf{Metric.} Per-window or per-event balanced F1 / AUROC; for TUSZ do
\emph{detection} (binary, every window labeled, whole-brain --- tidiest) before
\emph{typing} (multi-class, per-channel, imbalanced). Design the metric so a
\code{bckg}-only predictor does not look good (events/seizures are rare).

\smallskip
\textbf{Benchmark (optional).} Compare your event-level probe against a fine-tuned
\textbf{LaBraM}~\citep{LaBraM} foundation model on the same windows and splits.

\smallskip
\textbf{poc:} on the TUEV 6-class event task, frozen-JEPA $\sim$0.68 AUROC --- vs the
easy TUAB 0.877 of Track 1, this is where the recipe strains. TUAR and TUSZ are
preprocessed and ready as further \emph{independent} tasks --- each its own encoder,
pretrained and probed within that corpus (\emph{IID}).
\emph{Data:} \code{\$DSETS/Neuro/\{TUAR,TUEV,TUSZ\}\_PREPROCESSED} (each with its
\code{*\_README.md}).

\smallskip
\textbf{Stretch (cross-dataset transfer).} Beyond the IID setting, the three event
corpora are a natural transfer test: pretrain the JEPA on one, freeze it, then train a
\emph{fresh} probe on another's events (the label spaces differ, so only the encoder
transfers). Does an encoder trained on one kind of event generalize to another? That
generality is exactly what a foundation model claims.
\end{trackbox}

\begin{trackbox}{Track 3 --- Audio / speech representation JEPA \hfill \diffmed{}}
\textbf{Goal.} Self-supervised audio features, probed on keyword spotting (Speech
Commands) or environmental sound (ESC-50/UrbanSound). Treat log-mel spectrograms as
single-channel ``images'' for the image-JEPA masked/augmented setting, or treat mel
frames as a temporal sequence for the video-JEPA setting. \textbf{Files:} a
spectrogram dataset; reuse the ResNet or ViT encoder. \textbf{Metric:} linear-probe
accuracy on the downstream classification set. \textbf{Why JEPA:} audio is
continuous and noisy; latent prediction avoids modeling phase and background noise.

\smallskip
\textbf{poc:} Speech Commands (35-class, chance 2.86\%): linear-probe \textbf{72.6\%}
with log-mel spectrograms vs \textbf{33.2\%} with the raw 1D waveform; the
augmentation (VICReg) variant beat a masked-predictive one.
\emph{Data:} \code{/lustre/work/pdl17890/udl806719/datasets/speech\_commands\_v2} ~|~
\href{http://download.tensorflow.org/data/speech_commands_v0.02.tar.gz}{Speech Commands v0.02}
(Google; not on HF).
\end{trackbox}

\begin{trackbox}{Track 4 --- Wearable biosignals: ECG / IMU / PPG \hfill \diffeasy{}/\diffmed{}}
\textbf{Goal.} JEPA on accelerometer/gyroscope or ECG streams; probe on Human
Activity Recognition (UCI-HAR, PAMAP2) or arrhythmia (PhysioNet). These are
low-spatial multivariate time series, the ideal fit for the GRU predictor and
\code{[B, C, T, 1, 1]} convention. \textbf{Files:} a windowed time-series dataset;
1D-conv encoder; \code{RNNPredictor}; linear probe. \textbf{Metric:} probe accuracy
/ F1. The lowest-friction new-modality track, good for a first-time team.
\end{trackbox}

\begin{trackbox}{Track 5 --- Latent dynamics of physical fields (The Well) \hfill \diffmed{}/\diffhard{}}
\textbf{Goal.} The Well~\citep{the_well} is a 15\,TB collection of 16
physics-simulation datasets (fluid dynamics, magneto-hydrodynamics, astrophysics,
acoustics, active matter) stored as spatiotemporal fields, the continuous,
high-dimensional modality JEPAs were made for, with a ready-made PyTorch loader.
Train a video-JEPA-style model to predict the next field state in latent space, then
answer the open question: does latent prediction give more stable long-horizon
rollouts than the field-space neural-operator surrogates (FNO, U-Net) that The Well
ships as baselines? \textbf{Files:} wrap \code{the\_well.data.WellDataset} to emit
\code{[B, C, T, H, W]} clips; reuse the \code{ResUNet} spatial predictor and
\code{VCLoss}; add a small decoder from latent back to the field for scoring.
\textbf{Data.} Start at the small end of the 6.9\,GB--5.1\,TB range, a 2D set such
as \code{active\_matter} or \code{gray\_scott\_reaction\_diffusion}; do not pull the
multi-TB sets. \textbf{Metric:} VRMSE of the decoded autoregressive rollout vs the
FNO/U-Net baseline at several horizons. A clean JEPA-versus-surrogate comparison on
real physics (for \emph{observational} rather than simulated fields,
ERA5/WeatherBench is a harder variant).

\smallskip
\textbf{poc:} on \code{gray\_scott\_reaction\_diffusion}, the latent rollout beats a
persistence baseline \textbf{4.5$\times$} at horizon 6 and frozen latents separate
the 6 regimes at \textbf{0.92} accuracy (chance 0.167). But in decoded field-space
VRMSE the FNO/U-Net surrogates win at \emph{every} horizon $\le$10 --- a
latent$\to$field decoder floor ($\approx$0.19 VRMSE) caps the JEPA. Answer to the
open question: here the gain is representational, not pixel-accurate rollout.
\emph{Data:} \code{/lustre/work/pdl17890/udl806719/datasets/the\_well/gray\_scott\_reaction\_diffusion}
~|~ HF \href{https://huggingface.co/datasets/polymathic-ai/gray_scott_reaction_diffusion}{polymathic-ai/gray\_scott\_reaction\_diffusion}.
\end{trackbox}

\paragraph{Action-conditioned prediction: planning and forecasting (Tracks 6--9).}
The action-conditioned tracks below share one mapping onto the codebase. The
AC-video-JEPA (\cref{sec:acvideo}) is a latent dynamics model: given a state and an
\emph{action}, it predicts the next state's representation. The ``action'' is
whatever exogenous or control signal drives the system --- a task \emph{cue} for EEG
motor imagery (Track 6), or actuator commands, manipulated set-points, the operating
regime, known future covariates, or order-flow events for the time-series tracks
(7--9) --- and ``predict the future'' is forecasting the next state in latent space. Anomaly detection then falls out for
free: a fault or an attack is a \emph{violation of expectation}, a spike in
prediction energy against the \emph{realised} future (rather than against a
planning goal) that the action-conditioned dynamics cannot explain---the same
signal the intuitive-physics track uses on impossible videos. Concretely: fork
the AC-video-JEPA, feed it \code{(obs, action)} slices through a new
\code{TrajDataset} (\cref{sec:extension}), drop the planner, and read off either
the rollout forecast or the per-step energy. Two gotchas apply to all four: the
hardcoded \code{action\_dim} (\cref{sec:acvideo}) and mandatory per-channel
normalization (\cref{sec:extension}). With no control channel, the same recipe
degenerates to the video-JEPA forecasting setting (\cref{sec:video}).

\begin{trackbox}{Track 6 --- Mind in Motion: action-conditioned motor-imagery prediction (BCI~IV-2a) \hfill \diffmed{}/\diffhard{}}
\textbf{Usefulness in neuroscience.} Decoding \emph{intended} movement from EEG
(electroencephalography --- scalp recordings of the brain's electrical activity) is the
core of brain--computer interfaces that let paralysed patients drive a cursor or a
prosthetic.

\smallskip
\textbf{The idea (three parts, easy$\to$hard --- attempt all three).} BCI~IV-2a is
cue-based motor imagery: on each trial the subject imagines moving the left hand, right
hand, feet, or tongue. Every part starts from the \emph{same} JEPA encoder pretrained
self-supervised on these recordings (no labels); unlike the other tracks, here you are
expected to work through all three:
\begin{enumerate}[leftmargin=1.5em,itemsep=1pt,topsep=2pt]
  \item \emph{Within-subject probe (warm-up).} Freeze the encoder and linear-probe the 4
    classes per subject, fitting on session~\code{T} and testing on session~\code{E}
    (same person, different day) --- the original competition protocol.
  \item \emph{Cross-subject probe (hard).} The same probe, but \emph{leave-subjects-out}:
    fit on some subjects, test on held-out ones. Motor-imagery patterns are very
    person-specific, so this is the honest subject-disjoint test.
  \item \emph{Action-conditioned prediction (stretch).} The cue is an
    externally-imposed \emph{action}, so this is the one EEG corpus that maps onto the
    AC-video-JEPA (\cref{sec:acvideo}): encode the pre-cue baseline, condition the
    predictor on the cue, predict the post-cue imagery representation, and decode the cue
    by which action best reproduces the realised future (the inverse-dynamics head, or
    energy over the 4 actions). Encoder and predictor are trained \emph{jointly} with
    var/cov$+$IDM anti-collapse and no EMA target.
\end{enumerate}

\smallskip
\textbf{Data (preprocessed).} \code{BCICIV\_2a\_PREPROCESSED} --- 9 subjects $\times$
2 sessions, 288 balanced trials/session, 22 EEG channels at 250\,Hz, already
$0.5$--$100$\,Hz bandpass $+$ 50\,Hz notch (European data; the 3 EOG channels are
dropped). A per-trial cue$+$label \code{.csv} is provided for \emph{both} sessions;
\code{BCICIV\_2a\_README.md} gives the cue timing ($t{=}2$\,s), the imagery window
$[2,6]$\,s, and the labels.

\smallskip
\textbf{Files.} For the warm-up, a \code{Dataset} that loads the preprocessed EDF
(\code{MNE-Python}) and yields one epoch per sample in eb\_jepa's \code{[B, C, T, 1, 1]}
layout (\code{B} batch, \code{C} EEG channels, \code{T} time samples; the trailing
\code{1, 1} are the singleton spatial dims of the \code{[B, C, T, H, W]} convention), a
conv-patch / 1D-conv encoder in \file{architectures.py}, and a linear probe modeled on
\file{examples/image\_jepa/eval.py}. For the stretch, a \code{TrajDataset}
(\cref{sec:extension}) yielding \code{(pre-cue, cue, post-cue)} slices into the
AC-video-JEPA, the GRU \code{RNNPredictor}, and the IDM head as the cue decoder. Set the
hardcoded \code{action\_dim} to the 4-way cue.

\smallskip
\textbf{Metric.} 4-class accuracy / Cohen's $\kappa$ on held-out sessions (part~1) and
held-out subjects (part~2). For the action-conditioned part~3: does conditioning on the
\emph{correct} cue predict the realised future better than a wrong one, and what is the
IDM cue-decoding accuracy? Judge it on whether conditioning changes the prediction at
all, not on raw accuracy.

\smallskip
\textbf{Benchmark (optional).} For the probe parts, compare against a pretrained
\textbf{LaBraM}~\citep{LaBraM} foundation model on the same within-/cross-subject splits.

\smallskip
\textbf{Gotchas.} The EDF is continuous --- epoch each trial to its cue-locked imagery
window $[2,6]$\,s before probing. Tiny ($\sim$288 trials/subject); cross-subject is
genuinely hard (person-specific sensorimotor rhythms); EOG must never enter the
classifier (already dropped); 50\,Hz notch, not 60. Parts~1--2 are the safe deliverable;
part~3 is the ambitious one.
\emph{Data:} \code{\$DSETS/Neuro/BCICIV\_2a\_PREPROCESSED} (with
\code{BCICIV\_2a\_README.md}).
\end{trackbox}

% ===================== PASTE FROM HERE ========================================

\begin{trackbox}{Track 7 --- Scientific multivariate time-series forecasting \hfill \diffmed{}/\diffhard{}}
\textbf{Goal.} Pretrain a JEPA to predict the next window of a noisy scientific
multivariate series in latent space, then probe it on long-horizon forecasting or
regime classification. The research question: does anti-collapse latent prediction
learn more transferable, less overfit features than direct forecasting on
non-stationary signals?

\smallskip
\textbf{Data.} Use the \emph{de facto} long-horizon forecasting (LTSF) suite,
shipped with fixed train/val/test splits and $\sim$30 baselines in the
Time-Series-Library (\code{thuml/Time-Series-Library}) and introduced by
Informer~\citep{Informer} and Autoformer~\citep{Autoformer}: \textbf{Weather} (21
meteorological channels, 10-min), \textbf{ETT} (transformer oil temperature + 6
power-load channels, hourly / 15-min), \textbf{Electricity} (321 clients, hourly),
\textbf{Traffic} (862 road sensors, hourly), and \textbf{ILI} (weekly
influenza-like-illness counts). For breadth across scientific domains (solar,
temperature, hospital, nature), the Monash archive~\citep{MonashTSF} offers
$\sim$25 datasets in a single \code{.tsf} format (\code{Monash-University/monash\_tsf}
on Hugging Face). For spatiotemporal \emph{fields} rather than vector series, see
Track~5 (The Well) and ERA5/WeatherBench.

\smallskip
\textbf{Action-conditioned angle.} \textbf{ETT} is the cleanest fit: condition the
predictor on the six power-load channels (the ``action'') and forecast the oil
temperature---a genuinely controlled physical system (thermal response to load).
More generally, any dataset with \emph{known-future} exogenous covariates
(calendar, weather driving electricity / traffic) lets \code{q\_omega} encode the
covariate as the action; with none, drop to the video-JEPA forecasting setting.

\smallskip
\textbf{Files.} A windowed multivariate dataset wrapping the TSLib loaders (reuse
their splits verbatim for comparability); a 1D-conv or patch encoder in
\file{architectures.py}; the GRU \code{RNNPredictor} and the \code{[B, C, T, 1, 1]}
convention.

\smallskip
\textbf{Metric.} Forecasting MSE / MAE at the standard horizons
$\{96, 192, 336, 720\}$ (ILI: $\{24, 36, 48, 60\}$) against the published TSLib
numbers (DLinear, PatchTST, iTransformer, TimesNet), or a frozen-feature probe vs.\
a supervised baseline on the same split. Beating a strong linear baseline (DLinear)
is the bar worth clearing.

\smallskip
\textbf{poc.} Tested on \textbf{ETTh1} (\code{ailuntz/ETT-small}), TSLib 12/4/4 split,
input $L=336$, normalized MSE/MAE on test. A predictive JEPA (1D conv encoder +
\code{RNNPredictor} + EMA target + VICReg anti-collapse) was pretrained SSL on the
train-window inputs, then a frozen-feature ridge head forecast horizons
$\{96,192,336,720\}$, against a random-encoder floor, an end-to-end supervised head,
and an NLinear/DLinear baseline. \emph{The linear baseline dominates by $\sim$5$\times$}
(MSE \textbf{0.37 / 0.40 / 0.43 / 0.43}, matching the published DLinear numbers ---
confirming the split/normalization are correct), while frozen-JEPA $\approx$ random
$\approx$ supervised (MSE $\approx$1.2--2.9). Crucially the \emph{whole global-pooled
conv-encoder family} fails, not JEPA specifically: the supervised version of the same
backbone also loses to the linear map, reproducing the Zeng et al.\ (DLinear) lesson
that a pooled representation discards the level / last-value information ETT rewards;
latent prediction adds no transfer under this probe. \emph{Caveat:} the probe forecasts
from a single time-pooled vector --- a fairer JEPA-for-LTSF eval would keep a
per-patch / per-timestep representation (PatchTST-style) or fine-tune end-to-end, so this
is evidence against \emph{a global-pooled conv-JEPA representation} here, not against JEPA
forecasting in general ($n{=}1$ seed, ETTh1 only).
\emph{Data:} \code{/lustre/work/pdl17890/udl806719/datasets/LTSF/ETT-small}
~|~ HF \href{https://huggingface.co/datasets/ailuntz/ETT-small}{ailuntz/ETT-small}.
\end{trackbox}

\begin{trackbox}{Track 8 --- Financial time series: limit order books \& forecasting \hfill \diffmed{}/\diffhard{}}
\textbf{Goal.} Self-supervised features on high-frequency market data, probed on
mid-price-movement prediction. The question: can latent prediction extract a usable
signal from order-book dynamics that survives the brutal noise floor of financial
data?

\smallskip
\textbf{Data.} \textbf{FI-2010}~\citep{FI2010} is the standard public limit-order-book
(LOB) benchmark---often called the ``MNIST of high-frequency trading''---with
$\sim$4M samples from five Nasdaq Nordic stocks over ten days, the top ten LOB
levels (40 features: bid/ask price and volume), and 3-class mid-price-direction
labels (up / stationary / down) at horizons $k \in \{10, 20, 30, 50, 100\}$ ticks;
the standard protocol trains on the first seven days and tests on the last three.
For lower-frequency forecasting, the \textbf{Exchange-Rate} set (eight countries'
daily rates)~\citep{LSTNet} sits in the same LTSF suite as Track~7, and the Monash
archive~\citep{MonashTSF} adds further financial series.

\smallskip
\textbf{Action-conditioned angle.} The order book evolves under a stream of
events---placements, executions, cancellations. Treat that order flow (or
order-flow imbalance) as the ``action'' driving the book: condition the predictor
on the incoming event and forecast the next LOB state's representation, then read
the mid-price direction off the predicted latent.

\smallskip
\textbf{Files.} A LOB dataset (FI-2010 ships normalized; public PyTorch loaders
exist, e.g.\ the DeepLOB repositories); a temporal / 1D encoder; \code{RNNPredictor}.

\smallskip
\textbf{Metric.} 3-class mid-price F1 / accuracy at the standard horizons
(frozen probe vs.\ supervised), or forecasting MSE / MAE on Exchange-Rate.

\begin{gotchabox}
Markets are near-efficient and adversarial; mid-price direction is genuinely hard
and reported scores swing wildly with the label threshold and horizon. A clean
controlled comparison (JEPA features vs.\ supervised, with the collapse you fought)
is the win here---\emph{not} a number that looks like alpha. Treat any
``profitable'' result with deep suspicion and check for look-ahead leakage first.
\end{gotchabox}

\smallskip
\textbf{poc.} Tested on
\href{https://huggingface.co/datasets/thesven/fintime-decoder-dataset}{thesven/fintime-decoder-dataset}
(697 instruments, 2010--2025; windows of $87$ engineered features $\times\,64$ daily steps).
A predictive JEPA (1D conv encoder + \code{RNNPredictor} + EMA target + VICReg anti-collapse)
was pretrained SSL, frozen, and probed against a random-encoder floor and an end-to-end
supervised baseline. \emph{Negative result, exactly as the gotcha warns:} on the genuinely
hard targets (next-step direction / return) the frozen JEPA \emph{does not beat the random
encoder} (direction AUROC \textbf{0.58} vs random $0.58$ vs supervised $0.69$), and this
still holds on raw \emph{OHLCV-only} inputs (6 channels); returns are $\approx$chance
out-of-distribution for every method. The only non-trivial self-supervised signal comes
from \emph{VICReg, not latent prediction} (volatility-regime acc \textbf{0.94} vs random
$0.63$). \emph{Takeaway:} daily returns here are too close to efficient-market noise and the
engineered inputs are already informative, so this dataset cannot separate latent-prediction
from direct forecasting --- for a fair test of the track's hypothesis prefer a structured
series with learnable dynamics (the LTSF / Exchange-Rate suites of Track~7, or FI-2010 LOB
above).
\emph{Data:} \code{/lustre/work/pdl17890/udl806719/datasets/Finance/fintime\_prep}
~|~ HF \href{https://huggingface.co/datasets/thesven/fintime-decoder-dataset}{thesven/fintime-decoder-dataset}.
\end{trackbox}

\begin{trackbox}{Track 9 --- Action-conditioned anomaly prediction \& predictive maintenance \hfill \diffmed{}/\diffhard{}}
\textbf{The idea.} A JEPA trained only on \emph{healthy} operation learns the
action-conditioned dynamics of a system; at test time a fault, attack, or
degradation shows up as a spike in prediction energy---the model expected one
future and reality delivered another. This is forecasting-based anomaly detection,
and it is exactly the violation-of-expectation signal the intuitive-physics track
uses, now on real industrial data with a genuine control channel.

\smallskip
\textbf{Data (three flavours, each with an action channel).}
\begin{itemize}[leftmargin=1.4em,itemsep=2pt,topsep=2pt]
  \item \emph{Cyber-physical attacks.} \textbf{SWaT} and \textbf{WADI}~\citep{SWaT,WADI}
    are water treatment / distribution testbeds (iTrust, SUTD; request access) with
    both sensors \emph{and actuators}: SWaT has 51 channels over 7 normal + 4 attack
    days, WADI 123 channels over 14 + 2. Condition on the actuator commands (pump /
    valve states); the attacks are the anomalies.
  \item \emph{Process faults.} The \textbf{Tennessee Eastman Process}~\citep{TEP} is
    the canonical closed-loop chemical-plant benchmark: 41 measured + \emph{11
    manipulated} variables, 20 fault types, with a large multi-run simulation
    set~\citep{riethTEP} (also packaged by \code{AIRI-Institute/fddbenchmark}). The
    manipulated variables are the control actions, and the hard part is telling a
    fault from a normal closed-loop control response.
  \item \emph{Run-to-failure / RUL.} \textbf{NASA C-MAPSS}~\citep{CMAPSS} is
    \emph{the} prognostics benchmark: turbofan run-to-failure trajectories, 21
    sensors + \emph{3 operational settings} (altitude, Mach, throttle), four subsets
    FD001--FD004 of rising difficulty. Condition on the operational settings, track
    the drift of the latent, and regress remaining useful life.
\end{itemize}

\smallskip
\textbf{Methodology caveat.} The popular multivariate sets are flawed: Wu \&
Keogh~\citep{UCRanomaly} show the NASA telemetry sets (SMAP / MSL) contain trivial,
mislabeled, end-loaded anomalies, and even SWaT has many constant channels and
anomalies detectable from a single feature. For a clean, sound \emph{univariate}
evaluation use their \textbf{UCR Anomaly Archive} (250 curated signals)---it has no
action channel, so use it to validate the detector and SWaT / TEP / C-MAPSS to
exercise the action-conditioned story. Report event / range-based F1 and AUROC
(VUS-PR), not just point-adjusted F1, which inflates scores.

\smallskip
\textbf{Files.} A new \code{TrajDataset} yielding \code{(obs, action)} slices
(\cref{sec:extension}) plus an \code{env\_name} branch in
\file{datasets/utils.py:init\_data}; a 1D-conv encoder; \code{RNNPredictor}; the
AC-video-JEPA dynamics model (\cref{sec:acvideo}) with the planner removed and a
prediction-energy anomaly score. Watch the hardcoded \code{action\_dim}
(SWaT / WADI: dozens of actuators; TEP: 11; C-MAPSS: 3) and normalize every channel.

\smallskip
\textbf{Metric.} SWaT / WADI / TEP: event-based F1 and AUROC of the energy score,
healthy-only training. C-MAPSS: RUL RMSE and the NASA asymmetric score (which
penalises late predictions more) on FD001 first, then the harder FD004.
\end{trackbox}


\begin{trackbox}{Track 10 --- 3D point clouds / LiDAR \hfill \diffhard{}}
\textbf{Goal.} View-invariant JEPA on point clouds (ModelNet/ShapeNet): two
augmented samplings/rotations of an object are the two views (image-JEPA setting),
probed on shape classification. \textbf{Files:} a point-cloud dataset with geometric
augmentations; a PointNet-style encoder in \file{architectures.py}; reuse
\code{VICRegLoss}. \textbf{Metric:} linear-probe classification accuracy. Tests
whether the recipe transfers to an unordered, irregular modality.

\smallskip
\textbf{poc:} ModelNet40 (40-class, chance 2.5\%): linear-probe \textbf{80.9\%} with
aligned views, \textbf{64.7\%} under azimuth rotations, \textbf{35.3\%} under full
SO(3) --- accuracy drops monotonically as the demanded rotation invariance grows.
\emph{Data:} \code{/lustre/work/pdl17890/udl806719/datasets/modelnet40/modelnet40\_ply\_hdf5\_2048}
~|~ HF \href{https://huggingface.co/datasets/Msun/modelnet40}{Msun/modelnet40}.
\end{trackbox}

% ===========================================================================
\subsection{Group B: Push the vision and robotics examples}

\begin{trackbox}{Track 11 --- Learned cost / value for planning \hfill \diffmed{}/\diffhard{}}
\textbf{Goal.} Distance-in-latent-space is a crude planning cost. Replace it with a
learned cost or value function (TD-MPC style~\citep{TDMPC2,TD-MPC}) trained on the
world model's own rollouts, so the agent optimizes a quantity that correlates with
task success rather than raw representation distance. \textbf{Files:} a new
objective class in \file{eb\_jepa/planning.py} registered in
\code{objective\_name\_map}; a small value head trained alongside the JEPA.
\textbf{Metric:} Two Rooms success and planning compute vs the distance-cost
baseline, especially on the hardest non-monotonic episodes. Attacks a limitation
the paper explicitly calls out as open.

\smallskip
\textbf{poc:} on a frozen maze world model, the learned value reaches \textbf{37.5\%}
planning success vs \textbf{6.25\%} for the distance cost ($\sim$6$\times$) under
matched A$^*$ subgoals; greedy global planning is 0\% for both costs.
\emph{Data:} maze generated online (no dataset path / HF).
\end{trackbox}

\begin{trackbox}{Track 12 --- Hierarchical / multi-timescale world model \hfill \diffhard{}}
\textbf{Goal.} JEPAs here predict at a single temporal resolution; intelligent
planning wants several. Add a coarse predictor that models long-horizon
abstractions on top of the fine per-step predictor, and plan with the coarse level
to cut the search horizon. The paper names its modular encoder/predictor/regularizer
split as ``a natural starting point'' for exactly this, and no example ships it yet.
\textbf{Files:} a second predictor and a two-level unroll in a forked
\code{ac\_video\_jepa/main.py}; a coarse-level temporal-similarity term.
\textbf{Metric:} long-horizon planning success and wall-clock vs the flat baseline.
The deepest open robotics direction in the guide.

\smallskip
\textbf{poc:} on 21$\times$21 mazes, a two-level hierarchy solves them
\textbf{A$^*$-free at 65.6\% success / SPL 0.62} (matched-budget controls: 0\% for
greedy and 0\% for a random-walk baseline). Co-training the two levels did
\emph{not} beat the frozen single-level world model.
\emph{Data:} maze generated online (no dataset path / HF).
\end{trackbox}

\begin{trackbox}{Track 13 --- Stress-test the recipe under factors of variation \hfill \diffhard{}}
\textbf{Goal.} \texttt{eb\_jepa}'s Two Rooms is deliberately tiny. The open question
is whether the \emph{minimal} recipe (VC + sim + IDM, distance-cost MPPI) survives
realistic perturbations, and which term breaks first. Train an AC-video-JEPA-style
model and evaluate planning as you dial up controllable visual, geometric, and
physical factors of variation, using the environment suites in
\texttt{stable-worldmodel}~\citep{stableworldmodel} or the diverse planning
benchmarks in \texttt{jepa-wms}~\citep{terver2026JEPAWMs} rather than hand-building a
new environment. \textbf{Files:} adapt one of those suites' envs to the
\code{eb\_jepa} JEPA/planning stack (encoder + \code{JEPA.unroll} + a planner).
\textbf{Metric:} a success-vs-perturbation curve against the Two Rooms baseline, and
which regularizer term must be strengthened as the world gets harder. Ambitious but
high-signal: it tells you where the toy recipe actually ends.
\end{trackbox}

\begin{trackbox}{Track 14 --- Does intuitive physics emerge? \hfill \diffmed{}}
\textbf{Goal.} Recent work shows physical intuition emerges from self-supervised
video pretraining~\citep{garrido2025intuitivephysicsunderstandingemerges}. Test it
in miniature on the video-JEPA: feed it physically plausible vs impossible Moving
MNIST sequences (a digit teleporting, passing through a wall, or reversing
instantly) and ask whether prediction energy in latent space spikes on the
impossible ones, a violation-of-expectation signal. \textbf{Files:} a probe that
constructs paired plausible/impossible clips and compares \code{predcost} energy;
light additions to the video eval. \textbf{Metric:} the energy gap (impossible $>$
plausible) and how it grows over training. Exploratory, cheap, and genuinely
open at this scale.
\end{trackbox}

\begin{tipbox}
Whatever you pick, the deliverable that wins is a \emph{controlled comparison}: a
baseline and one change, three seeds each, one figure. ``We ported JEPA to EEG and
the linear probe hits 0.79 Balanced Accuracy, and here is the collapse we hit at
zero covariance weight and how we fixed it'' beats a pile of half-finished ideas.
\end{tipbox}
\section{Compute, Logistics, and Judging}
\label{sec:logistics}

\subsection{The compute budget}

We share \textbf{72 NVIDIA B200 GPUs} across \textbf{25 teams} for 24 hours. That
is about \textbf{2.9 GPUs per team} on average. The good news: every example in
this library is designed for \emph{single-GPU} training and finishes in a few
hours, so a team's natural unit of work is ``a few single-GPU jobs in parallel,''
not one giant distributed run. A three-seed sweep is three single-GPU jobs; a small
hyperparameter grid is a handful more.

\begin{notebox}
The numbers and the SLURM details below are the organizers' \emph{proposed}
policy. The exact \code{--partition}, \code{--account}, per-GPU memory, and any
fair-share quotas will be confirmed at kickoff and pinned in the event channel. If
a launch is rejected by the scheduler, it is almost certainly a partition/account
mismatch, ask in the channel rather than debugging it alone.
\end{notebox}

\paragraph{Fair-share rules of thumb.}
\begin{itemize}[leftmargin=1.4em,itemsep=1pt,topsep=2pt]
  \item Prefer many \emph{single-GPU} jobs over multi-GPU jobs. The examples do not
    need more than one GPU, and single-GPU jobs schedule faster and waste less.
  \item Cap your concurrency with \code{--array-parallelism}. A sane ceiling is
    \textbf{3 concurrent GPUs per team} during the day; if the queue is empty
    overnight, use more, but yield when teams are waiting.
  \item Kill dead runs. A collapsed run (flat variance, chance accuracy) will not
    recover; cancel it and free the GPU rather than letting it burn 12 hours.
  \item Smoke-test on \code{--single} with \code{optim.epochs} tiny before
    launching a sweep. A typo that crashes at epoch 0 across 12 array tasks wastes
    everyone's slots.
\end{itemize}

\paragraph{Disk and checkpoints.} Runs write to \code{\$EBJEPA\_CKPTS}. Checkpoints
add up fast across a sweep; keep only the best and the latest, and clean up smoke
tests. Datasets live in \code{\$EBJEPA\_DSETS}; download a shared copy once per
node rather than per team where possible (the organizers will pre-stage CIFAR-10,
Moving MNIST, and any agreed new-modality datasets).

\subsection{Notes on B200 vs the default H100 configs}

The shipped configs are tuned for H100s. B200s have substantially more memory and
strong \code{bfloat16} throughput, so the defaults will run but leave performance on
the table. Practical adjustments:

\begin{itemize}[leftmargin=1.4em,itemsep=1pt,topsep=2pt]
  \item \textbf{Increase batch size.} The image example uses \code{batch\_size=256}
    and the world model \code{384}; you can likely push these up on a B200. Scale
    the learning rate accordingly and re-check for collapse.
  \item \textbf{Keep \code{bfloat16}.} The configs already set \code{use\_amp: true}
    with \code{bfloat16} (image, world model) or \code{float16} (video). Prefer
    \code{bfloat16} on B200; if you see NaNs under \code{float16}, switch the video
    config to \code{bfloat16}.
  \item \textbf{Try \code{torch.compile}.} The world-model example already supports
    it; compilation pays off most on the longer temporal rollouts.
  \item \textbf{Do not over-scale.} A bigger batch is not free accuracy; for SSL it
    changes the regularization dynamics. Treat any batch-size change as an ablation
    with its own before/after, not a default win.
\end{itemize}

\begin{tipbox}
The bottleneck for these small models is often the data loader, not the GPU. If
GPU utilization is low, raise \code{data.num\_workers} and enable
\code{persistent\_workers}/\code{pin\_mem} before reaching for a bigger model. For a
new modality, a slow EDF/NetCDF reader can starve a B200 completely; pre-process to
a fast on-disk format (memmap, \code{.npy}, webdataset) once, up front.
\end{tipbox}

\subsection{Deliverables}

By the 24-hour mark each team submits:
\begin{enumerate}[leftmargin=1.6em,itemsep=1pt,topsep=2pt]
  \item \textbf{A repository}, submitted via \textbf{GitHub Classroom} (one repo per
    group): your fork/branch with all changes, runnable from a single documented
    command, plus the config(s) you used and a short report (hypothesis, setup,
    result with at least one figure or table, and one paragraph on what you learned
    about JEPAs --- especially any collapse you fought).
  \item \textbf{A 10-minute presentation} to the room (see \cref{sec:presentation}):
    your hypothesis, the controlled comparison, and one figure or live/recorded
    artifact (a probe-accuracy curve, a planning rollout GIF, a latent prediction).
\end{enumerate}

\subsection{The presentation}
\label{sec:presentation}

\emph{To be done.}

\subsection{How projects are judged}

\Cref{tab:rubric} is the rubric. Note what is \emph{not} on it: raw leaderboard
position. A negative result, cleanly shown (``latent prediction did not beat the
supervised baseline on this modality, and here is the evidence and our
hypothesis''), scores well. Confident hand-waving does not.

\begin{table}[h]
\centering
\caption{\textbf{Judging rubric.} Equal weight on understanding and execution;
creativity and clarity break ties.}
\label{tab:rubric}
\small
\begin{tabular}{lp{9.6cm}}
\toprule
Criterion & What we look for \\
\midrule
JEPA understanding   & Can the team explain encoder/predictor/regularizer, why collapse happens, and how their design avoids it? \\
Scientific rigor     & A clear hypothesis and a controlled before/after comparison; seeds and error bars where feasible. \\
Result               & A real, interpretable outcome (positive or negative) with a figure/number, not just ``it ran.'' \\
Creativity           & A modality or idea that genuinely stresses the recipe, or an elegant solution to a collapse/scaling problem. \\
Clarity              & A repo and presentation a non-expert teammate could follow. \\
\bottomrule
\end{tabular}
\end{table}

\subsection{Getting unstuck}

Read the gotcha boxes in \cref{sec:examples} first; most early problems are there.
For collapse, re-read \cref{sec:recipe} and watch the per-term losses. For anything
cluster-related, use the event channel. Mentors will run office hours on a posted
schedule. And remember the meta-advice: lock scope early, get a pipeline running in
hour five, and stop building in time to measure. Have a great 24 hours.

\bibliographystyle{iclr2026_conference}
\bibliography{main}

\end{document}
